\section{Conclusion}

This survey reviewed open-vocabulary and reasoning segmentation as a shift from fixed label sets toward language-conditioned, promptable, grounded, and reasoning-aware mask prediction. It first introduced the foundations and evaluation settings needed to compare this shift, then organized the literature into open-vocabulary segmentation, promptable and grounded segmentation, reasoning segmentation, and extended scenarios. Across these sections, the central trend is clear: visual segmentation is no longer only a dense classification problem, but an interface problem that connects language, spatial masks, domain knowledge, and user intent.

The reviewed methods show complementary strengths. CLIP-based approaches expand semantic coverage, SAM-style models improve promptable mask quality, grounded systems connect phrases and regions, and MLLM-based methods make implicit reasoning and conversational segmentation possible~\citep{radford2021learning,kirillov2023sam,ravi2024sam2,liang2023ovseg,xu2023san,xlaiLisaReasoningSegmentation2024,renPixellmPixelReasoning2024}. At the same time, the comparison shows that these gains are conditional on vocabulary design, prompt source, training data, foundation modules, metrics, and domain assumptions. The main value of this survey is therefore not to rank all methods under one number, but to provide a structured map for understanding what each method can prove and what remains uncontrolled.

Future progress will depend on transparent benchmarks, faithful reasoning-to-mask grounding, efficient hybrid systems, and domain-aware evaluation. These directions follow directly from the evidence reviewed above: as segmentation systems become more open and interactive, they must also become easier to verify, reproduce, and compare. A reliable next generation of open-vocabulary and reasoning segmentation should preserve the flexibility of language while making every predicted mask accountable to the visual evidence and task context that produced it.
